Los KPIs se definieron para decidir qu\'e combinaci\'on de t\'ecnica,
\textit{view} y perfil conviene recomendar. 

\begin{itemize}
    \item \textbf{Retorno total rebalanceado:}
    $R_{\text{total}} = \prod_t (1 + r_t) - 1$. Mide la ganancia acumulada
    fuera de muestra. Se usa para comparar el resultado final de cada cartera.

    \item \textbf{Retorno anualizado:}
    $R_{\text{anual}} = (1 + R_{\text{total}})^{1/T} - 1$. Normaliza el retorno
    por a\~no y permite comparar horizontes de distinta duraci\'on.

    \item \textbf{Volatilidad anualizada:}
    $\sigma_{\text{anual}} = \sigma_{\text{diaria}}\sqrt{252}$. Mide la
    variabilidad del retorno. Sirve para descartar alternativas demasiado
    inestables para el perfil del cliente.

    \item \textbf{Ratio de Sharpe rebalanceado:}
    $\text{Sharpe} = (R_{\text{anual}} - R_f)/\sigma_{\text{anual}}$, con
    $R_f = 0.02$. Resume retorno ajustado por riesgo. Es el principal KPI para
    comparar eficiencia financiera entre t\'ecnicas.

    \item \textbf{M\'aximo drawdown:} mayor ca\'ida porcentual desde un m\'aximo
    acumulado de riqueza. Eval\'ua la p\'erdida que el cliente habr\'ia tenido
    que tolerar, por lo que pesa especialmente en perfiles conservadores.

    \item \textbf{CVaR diario al 95\%:} p\'erdida promedio en el 5\% peor de
    d\'ias \cite{rockafellar2000}. Se usa para medir riesgo de cola y evitar
    carteras con p\'erdidas extremas demasiado severas.

    \item \textbf{Riqueza terminal esperada:} valor final promedio desde
    $C_0 = 1\,000$ USD en las 2\,000 simulaciones Monte Carlo por combinaci\'on. Traduce el
    desempe\~no financiero a una cifra monetaria para el cliente.

    \item \textbf{Probabilidad de ganancia:} proporci\'on de simulaciones en
    que la riqueza terminal supera $C_0$. Indica qu\'e tan probable es que el
    cliente termine el horizonte con utilidad.

    \item \textbf{Tasa de retiro:} proporci\'on de simulaciones en que el
    cliente abandona antes del final por superar su tolerancia de p\'erdida.
    Afecta la decisi\'on porque reduce la continuidad del cliente y las
    comisiones futuras.

    \item \textbf{Utilidad de FinPUC por saldo administrado:} ingreso esperado
    de la empresa, calculado como una comisi\'on mensual de 0,5\% sobre el saldo
    activo administrado de los clientes que permanecen en la plataforma. No
    incluye comisiones por transacci\'on ni por rebalanceo. Mide el atractivo de
    la alternativa para la empresa como consecuencia de la retenci\'on.

    \item \textbf{Clientes retenidos / saldo activo administrado:} n\'umero de
    clientes que permanecen activos al final del horizonte y saldo agregado bajo
    administraci\'on. Como la comisi\'on se cobra sobre el saldo activo, una
    mayor retenci\'on se traduce directamente en mayor utilidad para FinPUC.

    \item \textbf{Score P4:}
    \begin{equation}
    \label{eq:score_p4}
    \text{ScoreP4}_{v,m,i}
    =
    \bar{W}_{v,m,i}
    +
    \bar{U}_{v,m,i}
    -
    C_0 r_{v,m,i}.
    \end{equation}
    Para cada \textit{view} $v$, modelo $m$ y combinaci\'on
    $i=(s,p)$ de escenario hist\'orico y perfil de riesgo, $\bar{W}_{v,m,i}$
    corresponde a la riqueza terminal media del cliente, $\bar{U}_{v,m,i}$ a
    la utilidad media acumulada de FinPUC y $r_{v,m,i}$ a la tasa de abandono.
    Este \'indice resume riqueza del cliente, utilidad por saldo administrado y
    penalizaci\'on por retiro ($C_0$ por cada cliente que abandona). La utilidad
    de FinPUC entra como resultado \emph{ex post}; no forma parte del
    optimizador de portafolios.

    Para comparar cada \textit{view} Black-Litterman contra Markowitz base se
    calcula, en la misma combinaci\'on escenario--perfil, la mejora relativa:
    \begin{equation}
    \label{eq:delta_score_p4}
    \Delta_{v,i}
    =
    \frac{
    \text{ScoreP4}_{v,BL,i}
    -
    \text{ScoreP4}_{v,MK,i}
    }{
    |\text{ScoreP4}_{v,MK,i}|
    }.
    \end{equation}
    La m\'etrica final robusta de cada \textit{view} se obtiene como la mediana
    de estas mejoras relativas:
    \begin{equation}
    \label{eq:score_p4_mediana}
    \text{ScoreMediana}_{v}
    =
    \operatorname{mediana}_{i \in S \times P}
    \left(\Delta_{v,i}\right).
    \end{equation}
    La mediana mide el desempe\~no t\'ipico de la \textit{view} frente a
    Markowitz y reduce el efecto de observaciones extremas, especialmente
    perfiles muy arriesgados en escenarios con pandemia. Por ello, el informe
    reporta el Score P4 promedio como medida de potencial econ\'omico y el
    Score P4 por mediana como criterio principal de robustez comparativa.

    \item \textbf{Mejora de Sharpe $\Delta$Sharpe:} diferencia entre el Sharpe
    de un portafolio Black-Litterman calibrado y el Sharpe del Markowitz base
    equivalente, evaluados en el mismo horizonte de test. Un valor positivo
    indica que la calibraci\'on agrega valor sobre la estimaci\'on hist\'orica
    pura.

    \item \textbf{Turnover medio:} promedio de $\|w_t - w_{t-1}\|_1 / 2$ a lo
    largo del horizonte de evaluaci\'on. Mide la fracci\'on del portafolio que
    se rota en cada rebalanceo; valores altos implican mayores costos
    operacionales impl\'icitos.

    \item \textbf{N\'umero efectivo de activos:}
    $N_{\text{eff}} = 1 / \sum_i w_i^2$. Cuantifica la concentraci\'on del
    portafolio: $N_{\text{eff}} = N$ indica pesos uniformes, mientras que
    $N_{\text{eff}} \to 1$ indica concentraci\'on en pocos activos.

    \item \textbf{Concentraci\'on sectorial (HHI):}
    $\text{HHI} = \sum_{s} \left(\sum_{i \in s} w_i\right)^2$,
    donde $s$ indexa los sectores econ\'omicos. Valores cercanos a 1 indican
    concentraci\'on en un solo sector; valores bajos indican diversificaci\'on.
\end{itemize}
